\documentclass[11pt]{article}

\usepackage[margin=1in]{geometry}
\usepackage{amsmath,amssymb,mathtools}
\usepackage{booktabs}
\usepackage{enumitem}
\usepackage[hidelinks]{hyperref}

\newcommand{\R}{\mathbb{R}}
\newcommand{\Z}{\mathbb{Z}}
\newcommand{\ind}{\mathbb{I}}
\newcommand{\cvar}{\operatorname{CVaR}}

\title{Type-Indexed Distance-State CVaR Model}
\author{}
\date{}

\begin{document}
\maketitle

\section{Purpose}

This formulation jointly selects prepositioning locations and vehicle bases,
routes integer counts of interchangeable vehicles by type, allocates relief
resources, and minimizes the conditional value-at-risk (CVaR) of scenario
loss.  Individual vehicle labels are omitted.  Per-vehicle distance limits are
preserved by routing vehicle counts through a cumulative-distance-expanded
network.

The formulation below is the monolithic reference model.  A staged solution
method may solve only a subset of its decisions at a time, but it does not
change the definitions in this section.

\section{Sets and indices}

\begin{tabular}{ll}
\toprule
Symbol & Definition \\
\midrule
$N$ & Physical network nodes, indexed by $i,j,h$ \\
$N^P\subseteq N$ & Candidate prepositioning locations \\
$M$ & Transportation modes, indexed by $m$ \\
$K$ & Vehicle types, indexed by $k$ \\
$K_m$ & Vehicle types operating on mode $m$ \\
$R$ & Relief-resource types, indexed by $r$ \\
$\Omega$ & Disaster scenarios, indexed by $\omega$ \\
$A_m\subseteq N\times N$ & Directed physical arcs available to mode $m$ \\
$J_k\subseteq N^P$ & Eligible bases for vehicle type $k$ \\
\bottomrule
\end{tabular}

Let $m(k)$ denote the mode used by vehicle type $k$, and define
$A_k:=A_{m(k)}$.

\section{Parameters}

\begin{tabular}{ll}
\toprule
Symbol & Definition \\
\midrule
$F_k$ & Number of available vehicles of type $k$ \\
$D_k$ & Distance budget for each type-$k$ vehicle \\
$\psi_k$ & Turnaround time expressed as equivalent travel distance \\
$\operatorname{cap}_k$ & Payload capacity of one type-$k$ vehicle \\
$\ell_{ij}$ & Physical distance of arc $(i,j)$ \\
$c_{kij}$ & Unit resource-transport cost on arc $(i,j)$ \\
$w_r$ & Weight of one unit of resource $r$ \\
$d^\omega_{ir}$ & Demand for resource $r$ at node $i$ in scenario $\omega$ \\
$I_{ir}$ & Inventory available if candidate location $i$ is activated \\
$\rho$ & Releasable fraction of inventory \\
$\delta_{ir}$ & Unit unmet-demand penalty \\
$f_i$ & Fixed activation cost of candidate location $i$ \\
$B$ & Total activation budget \\
$P_{\max}$ & Maximum number of activated locations \\
$\pi^\omega$ & Probability of scenario $\omega$ \\
$\beta$ & CVaR confidence level \\
$\varepsilon$ & Small cost on vehicle movements \\
$a^\omega_i$ & Availability of node $i$ in scenario $\omega$ \\
$\sigma^\omega_i$ & Severity at node $i$ in scenario $\omega$ \\
$\gamma^\omega_m$ & Degradation sensitivity for mode $m$ in scenario $\omega$ \\
$T_{m,ij}$ & Nominal throughput of mode-$m$ arc $(i,j)$ \\
$\Theta_{i,m}$ & Baseline vehicle-handling capacity at node $i$ \\
$\Delta\Theta_{i,m}$ & Handling-capacity bonus if node $i$ is activated \\
\bottomrule
\end{tabular}

The scenario availability and degradation factors are
\begin{align}
a^\omega_i
&=
\begin{cases}
0, & \sigma^\omega_i\geq 1,\\
1, & \sigma^\omega_i<1,
\end{cases}
\label{eq:availability}\\
g^\omega_{m,ij}
&=
\max\left\{0,
1-\gamma^\omega_m
\frac{\max\{\sigma^\omega_i,\sigma^\omega_j\}}{5}
\right\},
\label{eq:arc-degradation}\\
\widehat C^\omega_{kij}
&=
\min\left\{
\operatorname{cap}_k,
T_{m(k),ij}g^\omega_{m(k),ij}
\right\}.
\label{eq:usable-capacity}
\end{align}

The residual node-handling capacity is
\begin{equation}
\widehat\Theta^\omega_{i,m}(p)
=
\max\left\{0,1-\gamma^\omega_m\frac{\sigma^\omega_i}{5}\right\}
\left(\Theta_{i,m}+\Delta\Theta_{i,m}p_i\right),
\label{eq:handling-capacity}
\end{equation}
where $p_i:=0$ for $i\notin N^P$ and
$\Delta\Theta_{i,m}:=0$ for $i\notin N^P$.

\section{Distance-state network}

Choose a number of distance buckets $H_k$ for each vehicle type.  In the
implementation, a common bucket count $H$ is normally used, such as $H=16$.
Define
\begin{equation}
\Delta_k=\frac{D_k}{H_k}
\label{eq:bucket-width}
\end{equation}
and the conservative bucket consumption of arc $(i,j)$ as
\begin{equation}
q_{kij}
=
\max\left\{1,
\left\lceil
\frac{\ell_{ij}+\psi_k}{\Delta_k}
\right\rceil
\right\}.
\label{eq:bucket-consumption}
\end{equation}

The admissible cumulative states are
\begin{equation}
\mathcal D_k=\{0,1,\ldots,H_k\}.
\end{equation}
The distance-expanded transition set is
\begin{equation}
\mathcal E_k^D
=
\left\{
(i,j,d,d'):
(i,j)\in A_k,\ d\in\mathcal D_k,\
d'=d+q_{kij}\leq H_k
\right\}.
\label{eq:expanded-arcs}
\end{equation}

For compactness, define the outgoing and incoming expanded transitions
\begin{align}
\delta_k^+(i,d)
&=
\{(i,j,d,d')\in\mathcal E_k^D\},\\
\delta_k^-(i,d)
&=
\{(h,i,d',d)\in\mathcal E_k^D\}.
\end{align}

For a physical arc $(i,j)$, define
\begin{equation}
\mathcal D_{kij}
=
\{d:(i,j,d,d+q_{kij})\in\mathcal E_k^D\}.
\end{equation}

\section{Decision variables}

First-stage decisions are
\begin{align}
p_i &\in\{0,1\}
&&\forall i\in N^P,
\label{eq:domain-p}\\
b_{kj} &\in\Z_+
&&\forall k\in K,\ j\in J_k,
\label{eq:domain-b}
\end{align}
where $b_{kj}$ is the number of type-$k$ vehicles based at $j$.

Scenario-dependent recourse decisions are
\begin{align}
n^\omega_{kijd} &\in\Z_+
&&\forall \omega,k,(i,j),d:
(i,j,d,d+q_{kij})\in\mathcal E_k^D,
\label{eq:domain-n}\\
\bar n^\omega_{kid} &\geq0
&&\forall\omega,k,i,d,
\label{eq:domain-nbar}\\
x^\omega_{kijr} &\geq0
&&\forall\omega,k,(i,j)\in A_k,r,
\label{eq:domain-x}\\
y^\omega_{ir} &\geq0
&&\forall\omega,i,r,
\label{eq:domain-y}\\
0\leq z^\omega_{ir} &\leq d^\omega_{ir}
&&\forall\omega,i,r,
\label{eq:domain-z}\\
\Lambda^\omega &\geq0
&&\forall\omega,
\label{eq:domain-loss}\\
\xi^\omega &\geq0
&&\forall\omega,
\label{eq:domain-xi}\\
\eta &\in\R.
\label{eq:domain-eta}
\end{align}

Here $n^\omega_{kijd}$ is the number of type-$k$ vehicles that traverse
physical arc $(i,j)$ after consuming $d$ distance buckets, and
$\bar n^\omega_{kid}$ is the number terminating at node $i$ in state $d$.

\section{Objective}

The model minimizes CVaR of scenario loss:
\begin{equation}
\min\quad
\eta+
\frac{1}{1-\beta}
\sum_{\omega\in\Omega}\pi^\omega\xi^\omega.
\label{eq:objective}
\end{equation}

\section{Constraints}

\subsection{Site activation and aggregate basing}

\begin{align}
\sum_{i\in N^P}p_i
&\leq P_{\max},
\label{eq:site-limit}\\
\sum_{i\in N^P}f_i p_i
&\leq B,
\label{eq:activation-budget}\\
\sum_{j\in J_k}b_{kj}
&=F_k
&&\forall k\in K,
\label{eq:fleet-assignment}\\
b_{kj}
&\leq F_k p_j
&&\forall k\in K,\ j\in J_k.
\label{eq:base-link}
\end{align}

\subsection{Distance-state vehicle conservation}

At every expanded node $(i,d)$,
\begin{align}
&\sum_{(i,j,d,d')\in\delta_k^+(i,d)}
n^\omega_{kijd}
+\bar n^\omega_{kid}
-\sum_{(h,i,d',d)\in\delta_k^-(i,d)}
n^\omega_{khid'}
\nonumber\\
&\hspace{35mm}=
\begin{cases}
a^\omega_i b_{ki},
& d=0,\ i\in J_k,\\
0,&\text{otherwise},
\end{cases}
&&\forall\omega\in\Omega,\ k\in K,\ i\in N,\ d\in\mathcal D_k.
\label{eq:vehicle-conservation}
\end{align}

Equation~\eqref{eq:vehicle-conservation} injects vehicle supply only at an
available selected base in distance state zero.  Every movement advances the
state by $q_{kij}\geq1$, and transitions beyond $H_k$ do not exist.

\subsection{Vehicle payload and degraded arc capacity}

For every physical vehicle arc,
\begin{equation}
\sum_{r\in R}w_r x^\omega_{kijr}
\leq
\widehat C^\omega_{kij}
\sum_{d\in\mathcal D_{kij}}n^\omega_{kijd}
\qquad
\forall\omega,k,(i,j)\in A_k.
\label{eq:vehicle-capacity}
\end{equation}

\subsection{Node handling capacity}

\begin{equation}
\sum_{k\in K_m}
\sum_{\substack{(h,i)\in A_k\\d\in\mathcal D_{khi}}}
n^\omega_{khid}
\leq
\widehat\Theta^\omega_{i,m}(p)
\qquad
\forall\omega\in\Omega,\ m\in M,\ i\in N.
\label{eq:node-handling}
\end{equation}

\subsection{Resource balance}

For each scenario, node, and resource,
\begin{align}
&\sum_{k\in K}
\sum_{(j,i)\in A_k}x^\omega_{kjir}
+\rho I_{ir}a^\omega_i p_i
+z^\omega_{ir}
\nonumber\\
&\hspace{25mm}=
d^\omega_{ir}
+y^\omega_{ir}
+\sum_{k\in K}
\sum_{(i,j)\in A_k}x^\omega_{kijr}
&&\forall\omega\in\Omega,\ i\in N,\ r\in R,
\label{eq:resource-balance}
\end{align}
where $I_{ir}p_i:=0$ for $i\notin N^P$.

\subsection{Scenario loss and CVaR}

Scenario loss is
\begin{align}
\Lambda^\omega
=
&\sum_{i\in N}\sum_{r\in R}
\delta_{ir}z^\omega_{ir}
\nonumber\\
&+
\sum_{k\in K}\sum_{(i,j)\in A_k}\sum_{r\in R}
c_{kij}x^\omega_{kijr}
\nonumber\\
&+
\varepsilon
\sum_{k\in K}\sum_{(i,j)\in A_k}
\sum_{d\in\mathcal D_{kij}}
n^\omega_{kijd}
&&\forall\omega\in\Omega.
\label{eq:scenario-loss}
\end{align}

The Rockafellar--Uryasev CVaR constraints are
\begin{equation}
\xi^\omega
\geq
\Lambda^\omega-\eta
\qquad
\forall\omega\in\Omega.
\label{eq:cvar-excess}
\end{equation}
Together with $\xi^\omega\geq0$ and objective~\eqref{eq:objective}, these
constraints imply
$\xi^\omega=\max\{0,\Lambda^\omega-\eta\}$ at an optimum.

\section{Per-vehicle distance guarantee}

For every expanded transition,
\begin{equation}
\ell_{ij}+\psi_k
\leq q_{kij}\Delta_k.
\end{equation}
For any individual path $P$ obtained by decomposing the integer vehicle flow,
\begin{align}
\sum_{(i,j)\in P}(\ell_{ij}+\psi_k)
&\leq
\Delta_k\sum_{(i,j)\in P}q_{kij}\\
&\leq
\Delta_k H_k\\
&=D_k.
\end{align}
Thus every decomposed vehicle path satisfies its distance budget; the limit is
not merely enforced on the fleet average.

\section{Relation to the staged algorithm}

The monolithic formulation optimizes all variables jointly over
$\Omega$.  The staged procedure uses the same detailed constraints but fixes
decisions sequentially:
\begin{enumerate}[label=\arabic*.]
\item choose $p$ and $b$ with either a reduced detailed model or a strategic
      proxy;
\item fix $p$ and $b$ and solve the detailed integer routing problem for each
      scenario;
\item fix $p$, $b$, and $n$, relax those fixed variables to continuous type,
      and solve the remaining all-scenario CVaR resource-allocation LP.
\end{enumerate}
The staged solution is feasible for the detailed formulation when every stage
is feasible, but it is not generally globally optimal because earlier fixed
decisions cannot be revised in the final CVaR solve.

\end{document}
